% Part: normal-modal-logic
% Chapter: frame-correspondence
% Section: equivalence-S5

\documentclass[../../../include/open-logic-section]{subfiles}

\begin{document}

\olfileid{nml}{frd}{es5}

\olsection{తుల్యతా సంబంధాలు మరియు \Log{S5}}

మోడల్ తర్కం \Log{S5}ను సార్వత్రిక చట్రాలన్నింటిలో
చెల్లుబాటు అయ్యే \tetoken{సూత్రాల}{!!{formula}s} సమితిగా
లక్షణీకరిస్తారు; అంటే ప్రతి లోకం నుంచి తనతో సహా
ప్రతి లోకం ప్రాప్యంగా ఉంటుంది. అలాంటి సందర్భంలో $\Box$
అనివార్యతకు, $\Diamond$ సాధ్యతకు అనురూపం: ప్రతి
లోకంలోనూ $!A$ సత్యమైతే $\Box !A$ సత్యం;
ఏదో ఒక లోకంలో $!A$ సత్యమైతే $\Diamond !A$ సత్యం.
\Log{S5}ను అన్ని స్వావర్తన, సౌష్ఠవ, సంక్రామక
చట్రాల్లో, అంటే \emph{తుల్యతా సంబంధాలు} గల చట్రాల్లో,
చెల్లుబాటు అయ్యే \tetoken{సూత్రాల}{!!{formula}s} సమితిగా
కూడా లక్షణీకరించవచ్చు.

\begin{defn}
  $W$పై ద్విస్థాన సంబంధం $R$ స్వావర్తనమూ,
  సౌష్ఠవమూ, సంక్రామకమూ అయితే, అప్పుడు మాత్రమే
  దాన్ని \emph{తుల్యతా సంబంధం} అంటాం. $W$పై సంబంధం
  $R$ అన్ని $u,v \in W$లకు $Ruv$ అయితే,
  అప్పుడు మాత్రమే దాన్ని \emph{సార్వత్రిక సంబంధం} అంటాం.
\end{defn}

\Ax{T}, \Ax{B}, \Ax{4} వరుసగా స్వావర్తన,
సౌష్ఠవ, సంక్రామక చట్రాలను లక్షణీకరిస్తాయి కాబట్టి,
ప్రాప్యత సంబంధం తుల్యతా సంబంధమైన చట్రాలు ఖచ్చితంగా
ఈ మూడు \tetoken{సూత్రాలూ}{!!{formula}s} చెల్లుబాటు అయ్యే
చట్రాలే. తుల్యతా సంబంధాలను ఇతర
\tetoken{సూత్రాల}{!!{formula}s} కలయికలతో కూడా
లక్షణీకరించవచ్చు: వాటి నిర్వచనానికి వాడిన షరతులు,
$R$పై మనకు తెలిసిన ఇతర షరతుల కలయికలకు సమానమైనవి.

\begin{prop}\ollabel{prop:equivalences}
  కింది షరతులు సమానార్థకాలు:
  \begin{enumerate}
  \item $R$ తుల్యతా సంబంధం;
  \item $R$ స్వావర్తనమూ, యూక్లిడియన్ కూడా;
  \item $R$ సీరియల్, సౌష్ఠవ, యూక్లిడియన్;
  \item $R$ సీరియల్, సౌష్ఠవ, సంక్రామక.
  \end{enumerate}
\end{prop}

\begin{proof}
  అభ్యాసం.
\end{proof}

\begin{prob}
  కింది విషయాలను చూపి
  \olref[nml][frd][es5]{prop:equivalences}ను నిరూపించండి:
  \begin{enumerate}
  \item $R$ సౌష్ఠవమూ, సంక్రామకమూ అయితే అది యూక్లిడియన్.
  \item $R$ స్వావర్తనమైతే అది సీరియల్.
  \item $R$ స్వావర్తనమూ, యూక్లిడియన్ అయితే అది సౌష్ఠవం.
  \item $R$ సౌష్ఠవమూ, యూక్లిడియన్ అయితే అది సంక్రామకం.
  \item $R$ సీరియల్, సౌష్ఠవ, సంక్రామక అయితే అది స్వావర్తనం.
  \end{enumerate}
  ఈ విషయాలు నాలుగు షరతుల సమానార్థకతను నిరూపించడానికి
  ఎందుకు సరిపోతాయో వివరించండి.
\end{prob}

\olref{prop:equivalences},
\olref[prf][prs]{prop:S5}కు అర్థపరమైన అనురూప ఫలితం:
ప్రాప్యత సంబంధం~$R$ తుల్యతా సంబంధమైన చట్రాల మోడల్
తర్కానికి, సంప్రదాయంగా \Log{S5} అని పిలిచే తర్కానికి,
అది సమానమైన మరో లక్షణీకరణను ఇస్తుంది.

సార్వత్రిక సంబంధాలకూ తుల్యతా సంబంధాలకూ సంబంధం ఏమిటి?
ప్రతి సార్వత్రిక సంబంధమూ తుల్యతా సంబంధమే అయినా,
ప్రతి తుల్యతా సంబంధమూ సార్వత్రికం కాదని స్పష్టం.
అయినా సార్వత్రిక సంబంధాలన్నింటిలో చెల్లుబాటు అయ్యే
\tetoken{సూత్రాల}{!!{formula}s} సమితి, తుల్యతా
సంబంధాలన్నింటిలో చెల్లుబాటు అయ్యే సూత్రాల
సమితితో ఖచ్చితంగా ఒకటే.

\begin{prop}
  $R$ ఒక తుల్యతా సంబంధం అనుకుందాం. ప్రతి $w \in W$కు,
  $w$ యొక్క \emph{తుల్యతా వర్గాన్ని} $[w] = \{w'\in W :
  Rww'\}$ అనే సమితిగా నిర్వచిద్దాం. అప్పుడు:
  \begin{enumerate}
  \item $w \in [w]$;
  \item ప్రతి తుల్యతా వర్గం $[w]$పై $R$ సార్వత్రికం;
  \item తుల్యతా వర్గాల సముదాయం $W$ను పరస్పరం
    వియుక్తమై, కలిపి మొత్తాన్ని కప్పే ఉపసమితులుగా
    విభజిస్తుంది.
  \end{enumerate}
\end{prop}

\begin{prop}\ollabel{prop:S5=univ}
  తుల్యతా సంబంధం $R$ గల అన్ని చట్రాలు
  $\mModel{F} =\tuple{W,R}$లో
  \tetoken{ఒక సూత్రం}{!!^a{formula}}~$!A$ చెల్లుబాటు
  అయ్యేది, సార్వత్రిక సంబంధం $R$ గల అన్ని చట్రాలు
  $\mModel{F} =\tuple{W,R}$లో అది చెల్లుబాటు
  అయినప్పుడే, అప్పుడే. కనుక సార్వత్రిక చట్రాల
  తర్కం \Log{S5}యే.
\end{prop}

\begin{proof}
  $W$పై సార్వత్రిక సంబంధం~$R$ తుల్యతా సంబంధమని
  వెంటనే తనిఖీ చేయవచ్చు. కనుక $R$ తుల్యతా సంబంధం గల
  అన్ని చట్రాల్లో $!A$ చెల్లుబాటు అయితే,
  సార్వత్రిక చట్రాలన్నింటిలోనూ చెల్లుబాటు అవుతుంది.
  విలోమ దిశకు వ్యతిరేక ప్రతిజ్ఞను నిరూపిద్దాం:
  $W$పై $R$ తుల్యతా సంబంధమైన చట్రం
  $\tuple{W,R}$పై ఆధారపడే నమూనా
  $\mModel{M} = \tuple{W, R, V}$లోని ఒక లోకం $w$
  వద్ద \tetoken{ఒక సూత్రం}{!!a{formula}}~$!B$
  విఫలమవుతుందనుకుందాం. అంటే $\mSat/{M}{!B}[w]$.
  ఇప్పుడు $\mModel{M}' = \tuple{W', R', V'}$ నమూనాను
  ఇలా నిర్వచిద్దాం:
  \begin{enumerate}
  \item $W' = [w]$;
  \item $W'$పై $R'$ సార్వత్రికం;
  \item $V'(p) = V(p) \cap W'$.
  \end{enumerate}
  (కాబట్టి $\mModel{M}'$లోని లోకాల సమితి $W'$ను
  \olref{fig:partition}లో మసక భాగం సూచిస్తుంది.)
  $W'$పై $R$, $R'$ ఒకే సంబంధమని సులభంగా చూడవచ్చు.
  తరువాత \tetoken{సూత్రాలపై}{!!{formula}s} ఆగమనంతో,
  ప్రతి $!A$కు, అన్ని $w' \in W'$ల వద్ద
  $\mSat{M'}{!A}[w']$ అయినప్పుడే, అప్పుడే
  $\mSat{M}{!A}[w']$ అని చూపవచ్చు ($W'
  \subseteq W$ కాబట్టి ఈ పోలిక అర్థవంతమే).
  ముఖ్యంగా $\mSat/{M'}{!B}[w]$; కాబట్టి
  సార్వత్రిక చట్రంపై ఆధారపడే ఒక నమూనాలో $!B$
  విఫలమవుతుంది.
\end{proof}

\begin{figure}[t]
  \centering
  \begin{tikzpicture}[node distance=2cm, auto, thick]
    \clip [rounded corners] (0,0) -- (6,0) -- (6,4) -- (0,4) --  cycle;
    \begin{scope}
      \clip (0,0) -- (2,0)  .. controls (1.5,1.5) and (4.5,1.5)
      .. (4,4) -- (0,4) -- (0,0) -- cycle;
      \filldraw[fill=gray!40] (0,4) -- (0,2) .. controls (1.5,1.5)
      and (4.5,1.5) .. (4.5,0) -- (6,0) -- (6,4) -- (0,4) --  cycle;
    \end{scope}
    \draw [name path=line1] (2,0) .. controls (1.5,1.5) and (4.5,1.5) .. (4,4);
    \draw [name path=line2] (0,2)  .. controls (1.5,1.5) and (4.5,1.5) .. (4.5,0);
    \path [name intersections={of=line1 and line2}];
    \draw [rounded corners, use as bounding box, very thick] (0,0)
    -- (6,0) -- (6,4) -- (0,4) --  cycle;
    \node at (2,3) {$[w]$} ;
    \node at (1,1) {$[u]$} ;
    \node at (3,0.75) {$[v]$} ;
    \node at (4.75,2) {$[z]$} ;
  \end{tikzpicture}
  \caption{$W$ను తుల్యతా వర్గాలుగా విభజించడం.}
  \ollabel{fig:partition}
\end{figure}

\end{document}
